%========================================
% MatinBook — Stage 14: Complete Book Test
% Version: v1.1
% Purpose:
%   Full integration test of all MatinBook v1.1 features:
%     1. Full book structure (cover, frontmatter, mainmatter,
%        backmatter, back cover)
%     2. All box families (matinbox, matin-outline)
%     3. Math environments and operators
%     4. Code listings (minted) with Python
%     5. Algorithms with Persian captions
%     6. Graphics (TikZ + PGFPlots)
%     7. Bibliography (biblatex + biber)
%     8. Index (xindy with Persian sorting)
%     9. Cross-references (\cref)
%    10. Margin notes (\mnote)
%
% Compilation:
%   xelatex -shell-escape stage14-book.tex
%   biber stage14-book
%   xelatex -shell-escape stage14-book.tex
%   xelatex -shell-escape stage14-book.tex
%   (full cycle required for bibliography + index + cover)
%
% Expected outcome:
%   A complete, professionally typeset book demonstrating
%   that all MatinBook modules work together coherently.
%========================================

%------------------------------------------------------------------------
% Search Path (for tests directory)
%------------------------------------------------------------------------
\makeatletter
\def\input@path{{../}}
\makeatother

%------------------------------------------------------------------------
% Document Class
%------------------------------------------------------------------------
\documentclass{matinbook}

%------------------------------------------------------------------------
% Bibliography Resource
%------------------------------------------------------------------------
\addbibresource{../references.bib}

%------------------------------------------------------------------------
% Book Metadata
%------------------------------------------------------------------------
\title{راهنمای جامع برنامه‌نویسی و الگوریتم}
\author{تیم توسعه MatinBook}
\date{مهر ۱۴۰۵}

\booktitle{راهنمای جامع برنامه‌نویسی و الگوریتم}

\renewcommand{\repository}{github.com/matinmahmoudi-cs/matinbook}

%========================================================================
% DOCUMENT
%========================================================================
\begin{document}

%------------------------------------------------------------------------
% FRONT COVER
%------------------------------------------------------------------------
\makecover
    {راهنمای جامع برنامه‌نویسی و الگوریتم}
    {از مفاهیم پایه تا پیشرفته}
    {تیم توسعه MatinBook}
    {مهر ۱۴۰۵}

%------------------------------------------------------------------------
% FRONT MATTER (symmetric margins)
%------------------------------------------------------------------------
\frontmatter
\frontmattergeometry

% Title page
\maketitle

% Copyright / Colophon
\thispagestyle{empty}
\vfill
\begin{center}
    {\large\textbf{حقوق نشر}}

    \vspace{1cm}

    تمامی حقوق این کتاب برای تیم \lr{MatinBook} محفوظ است.
    هرگونه کپی‌برداری، تکثیر و انتشار بدون اجازه‌ی کتبی ممنوع است.

    \vspace{1cm}

    نسخه \matinbookversion\ — \matinbookdate

    \vspace{1cm}

    \lr{Copyright \copyright\ 2026 MatinBook Project} \\
    \lr{Licensed under the MIT License}
\end{center}
\clearpage

% Dedication
\thispagestyle{empty}
\vfill
\begin{flushright}
    {\Large\itshape
        تقدیم به تمام جویندگان دانش\\
        که هیچ‌گاه از یادگیری دست نمی‌کشند.
    }
\end{flushright}
\clearpage

% Preface
\chapter*{پیشگفتار}
\addcontentsline{toc}{chapter}{پیشگفتار}

\section*{درباره‌ی این کتاب}

این کتاب حاصل ماه‌ها تلاش تیم \lr{MatinBook} برای ایجاد
یک منبع آموزشی کامل در زمینه‌ی برنامه‌نویسی و الگوریتم
است. هدف ما ارائه‌ی مطالبی دقیق، روان و کاربردی با
استفاده از ابزارهای حروف‌چینی مدرن فارسی است.

\section*{ساختار کتاب}

کتاب در سه بخش اصلی سازماندهی شده است:

\begin{enumerate}
    \item \textbf{بخش اول — مبانی برنامه‌نویسی:}
    آشنایی با پایتون، ساختارهای داده و کنترل جریان.
    \item \textbf{بخش دوم — الگوریتم‌ها:}
    الگوریتم‌های جستجو، مرتب‌سازی و تحلیل پیچیدگی.
    \item \textbf{بخش سوم — مفاهیم پیشرفته:}
    برنامه‌نویسی شیءگرا، ساختمان داده‌های پیشرفته و
    الگوریتم‌های گراف.
\end{enumerate}

\section*{نحوه‌ی استفاده}

هر فصل شامل توضیحات نظری، کدهای نمونه، نمودارهای
آموزشی و تمرین‌های عملی است. توصیه می‌شود خواننده
علاوه بر مطالعه‌ی متن، کدها را اجرا و تمرین‌ها را
حل کند.

\vspace{1cm}
\begin{flushleft}
    \textbf{تیم توسعه MatinBook} \\
    مهر ۱۴۰۵
\end{flushleft}

\clearpage

% Table of Contents
\tableofcontents
\clearpage

% List of Figures
\listoffigures
\clearpage

% List of Tables
\listoftables
\clearpage

% List of Algorithms
\listofalgorithms
\clearpage

%------------------------------------------------------------------------
% MAIN MATTER (wide margin for notes)
%------------------------------------------------------------------------
\mainmatter
\mainmattergeometry

%========================================================================
% PART 1 — PROGRAMMING BASICS
%========================================================================
\part{مبانی برنامه‌نویسی}

%========================================================================
% CHAPTER 1 — PYTHON FUNDAMENTALS
%========================================================================
\chapter{مبانی پایتون}
\label{chap:python}

\section{معرفی زبان پایتون}
\label{sec:python-intro}

پایتون\index{پایتون} یک زبان برنامه‌نویسی سطح بالا،
تفسیری و همه‌منظوره است که توسط خیدو فان
روسوم\index{فان روسوم، خیدو} در سال ۱۹۹۱ میلادی
ساخته شد. این زبان بر خوانایی کد و سادگی نحو تأکید
دارد و امروزه در حوزه‌های مختلفی از جمله علم
داده\index{علم داده}، هوش مصنوعی\index{هوش مصنوعی}،
توسعه‌ی وب\index{توسعه وب} و خودکارسازی\index{خودکارسازی}
استفاده می‌شود.

\mnote{پایتون در سال ۲۰۲۴ محبوب‌ترین زبان برنامه‌نویسی جهان بوده است.}

\begin{definition}[پایتون]
    پایتون یک زبان برنامه‌نویسی شیءگرا، تفسیری و با
    نوع‌دهی پویا (\lr{dynamic typing}) است که فلسفه‌ی
    طراحی آن بر خوانایی کد و کاهش پیچیدگی نحو استوار
    است.
    \label{def:python}
\end{definition}

\section{انواع داده و متغیرها}
\label{sec:python-datatypes}

انواع داده‌ی پایه در پایتون در \cref{code:data-types}
نمایش داده شده است:

\begin{latin}
\begin{minted}{python}
# Basic data types in Python
name: str = "Ali"                    # String
age: int = 25                        # Integer
height: float = 1.75                 # Float
is_student: bool = True              # Boolean

# Collection types
grades: list = [18, 17, 20, 19]     # List
point: tuple = (10, 20)              # Tuple
info: dict = {"name": "Ali"}         # Dictionary
tags: set = {"python", "coding"}     # Set

# Print all values
print(f"Name: {name}, Age: {age}")
print(f"Height: {height}, Student: {is_student}")
\end{minted}
\end{latin}
\codecaption{انواع داده پایه در پایتون}
\label{code:data-types}

\section{ساختارهای کنترلی}
\label{sec:python-control}

\begin{definition}[ساختار شرطی]
    ساختار شرطی\index{ساختار شرطی} به برنامه اجازه
    می‌دهد بر اساس شرایط مختلف، مسیرهای متفاوتی را
    اجرا کند.
    \label{def:conditional}
\end{definition}

نمونه‌ای از ساختار شرطی در \cref{code:conditional}:

\begin{latin}
\begin{minted}{python}
def classify_number(n: int) -> str:
    """
    Classify a number as positive, negative, or zero.

    Args:
        n: An integer to classify.

    Returns:
        A string describing the number.
    """
    if n > 0:
        return "Positive"
    elif n < 0:
        return "Negative"
    else:
        return "Zero"

# Test with various inputs
test_numbers = [5, -3, 0, 42, -17]
for num in test_numbers:
    result = classify_number(num)
    print(f"{num:4d} -> {result}")
\end{minted}
\end{latin}
\codecaption{ساختار شرطی در پایتون}
\label{code:conditional}

\section{حلقه‌ها و تکرار}
\label{sec:python-loops}

انواع حلقه‌ها در \cref{code:loops}:

\begin{latin}
\begin{minted}{python}
# For loop with range
print("Counting with for loop:")
for i in range(5):
    print(f"  Iteration {i}")

# While loop
print("\nCounting with while loop:")
count = 0
while count < 5:
    print(f"  Count: {count}")
    count += 1

# List comprehension
squares = [x**2 for x in range(10)]
print(f"\nSquares: {squares}")
\end{minted}
\end{latin}
\codecaption{انواع حلقه‌ها در پایتون}
\label{code:loops}

%========================================================================
% CHAPTER 2 — FUNCTIONS AND MODULES
%========================================================================
\chapter{توابع و ماژول‌ها}
\label{chap:functions}

\section{تعریف تابع}
\label{sec:functions-definition}

تابع\index{تابع} بلوکی از کد است که می‌تواند ورودی
دریافت کند، عملیاتی روی آن انجام دهد و خروجی تولید
کند.

\begin{definition}[تابع در پایتون]
    تابع با کلمه‌ی کلیدی \inlcode{def} تعریف می‌شود
    و می‌تواند پارامترهای ورودی، مستندات
    (\lr{docstring}) و مقدار بازگشتی داشته باشد.
    \label{def:function}
\end{definition}

نمونه‌ی تابع بازگشتی در \cref{code:factorial}:

\begin{latin}
\begin{minted}{python}
def factorial(n: int) -> int:
    """
    Calculate the factorial of n recursively.

    Args:
        n: A non-negative integer.

    Returns:
        n! = n * (n-1) * ... * 1

    Raises:
        ValueError: If n is negative.
    """
    if n < 0:
        raise ValueError("Factorial is not defined for negative numbers")
    if n <= 1:
        return 1
    return n * factorial(n - 1)


# Calculate and display factorials
for i in range(10):
    print(f"{i}! = {factorial(i)}")
\end{minted}
\end{latin}
\codecaption{تابع فاکتوریل بازگشتی}
\label{code:factorial}

\begin{example}
    خروجی تابع \cref{code:factorial}:
    \begin{align*}
        0! &= 1 \\
        5! &= 5 \times 4 \times 3 \times 2 \times 1 = 120 \\
        10! &= 3628800
    \end{align*}
    \label{ex:factorial-values}
\end{example}

\section{آرگومان‌های تابع}
\label{sec:functions-arguments}

آرگومان‌های متغیر در \cref{code:arguments}:

\begin{latin}
\begin{minted}{python}
def describe_person(
    name: str,
    age: int = 0,
    *hobbies: str,
    **details
) -> str:
    """
    Create a description of a person.
    """
    parts = [f"Name: {name}"]

    if age > 0:
        parts.append(f"Age: {age}")

    if hobbies:
        parts.append(f"Hobbies: {', '.join(hobbies)}")

    for key, value in details.items():
        parts.append(f"{key}: {value}")

    return "\n".join(parts)


# Test the function
print(describe_person("Ali", 25, "coding", "reading"))
print("---")
print(describe_person("Sara", job="Engineer", city="Tehran"))
\end{minted}
\end{latin}
\codecaption{آرگومان‌های متغیر در پایتون}
\label{code:arguments}

%========================================================================
% PART 2 — ALGORITHMS
%========================================================================
\part{الگوریتم‌ها}

%========================================================================
% CHAPTER 3 — SEARCH ALGORITHMS
%========================================================================
\chapter{الگوریتم‌های جستجو}
\label{chap:search}

\section{جستجوی خطی}
\label{sec:search-linear}

الگوریتم جستجوی خطی\index{الگوریتم!جستجوی خطی}
ساده‌ترین روش جستجو است. در این روش، تک‌تک عناصر
آرایه بررسی می‌شوند تا عنصر مورد نظر پیدا شود
(\cref{alg:linear}).

\begin{latin}
\begin{algorithm}
\caption{جستجوی خطی}
\label{alg:linear}
\begin{algorithmic}[1]
    \Require $A[1..n]$: array of $n$ elements
    \Require $x$: target value to find
    \Ensure Index of $x$ in $A$, or $-1$ if not found

    \For{$i \gets 1$ \textbf{to} $n$}
        \If{$A[i] = x$}
            \State \Return $i$
        \EndIf
    \EndFor
    \State \Return $-1$
\end{algorithmic}
\end{algorithm}
\end{latin}

\begin{theorem}[پیچیدگی جستجوی خطی]
    الگوریتم \cref{alg:linear} در بدترین حالت
    \lr{$O(n)$} مقایسه انجام می‌دهد.
    \label{thm:linear-complexity}
\end{theorem}

\section{جستجوی دودویی}
\label{sec:search-binary}

برای آرایه‌های \textbf{مرتب}، جستجوی
دودویی\index{الگوریتم!جستجوی دودویی} بسیار
کارآمدتر است (\cref{alg:binary}).

\begin{latin}
\begin{algorithm}
\caption{جستجوی دودویی}
\label{alg:binary}
\begin{algorithmic}[1]
    \Require $A[1..n]$: sorted array (ascending)
    \Require $x$: target value to find
    \Ensure Index of $x$ in $A$, or $-1$ if not found

    \State $left \gets 1$, $right \gets n$
    \While{$left \leq right$}
        \State $mid \gets \lfloor (left + right) / 2 \rfloor$
        \If{$A[mid] = x$}
            \State \Return $mid$
        \ElsIf{$A[mid] < x$}
            \State $left \gets mid + 1$
        \Else
            \State $right \gets mid - 1$
        \EndIf
    \EndWhile
    \State \Return $-1$
\end{algorithmic}
\end{algorithm}
\end{latin}

نمودار \cref{fig:search-comparison} رشد پیچیدگی دو
الگوریتم را مقایسه می‌کند:

\begin{figure}[h]
\centering
\begin{latin}
\begin{tikzpicture}
  \begin{axis}[
    width=\textwidth, height=6.5cm,
    xlabel={Input Size ($n$)},
    ylabel={Number of Comparisons},
    title={Search Algorithm Comparison},
    grid=major,
    legend pos=north west,
    domain=0:100,
    samples=100
  ]
    \addplot[matinred, thick] {x};
    \addlegendentry{Linear Search $O(n)$}
    \addplot[matinblue, thick] {log2(x)};
    \addlegendentry{Binary Search $O(\log n)$}
  \end{axis}
\end{tikzpicture}
\end{latin}
\caption{مقایسه رشد پیچیدگی الگوریتم‌های جستجو}
\label{fig:search-comparison}
\end{figure}

\begin{remark}
    جستجوی دودویی (\cref{alg:binary}) فقط روی
    آرایه‌های \textbf{مرتب} کار می‌کند. اگر آرایه
    مرتب نباشد، ابتدا باید آن را مرتب کرد که خود
    هزینه‌ی \lr{$O(n \log n)$} دارد.
    \label{rem:binary-requirement}
\end{remark}

%========================================================================
% CHAPTER 4 — SORTING ALGORITHMS
%========================================================================
\chapter{الگوریتم‌های مرتب‌سازی}
\label{chap:sorting}

\section{مرتب‌سازی سریع}
\label{sec:sorting-quicksort}

\lr{Quick Sort}\index{الگوریتم!مرتب‌سازی سریع} از
الگوریتم‌های تقسیم و حل\index{تقسیم و حل} است که
در عمل بسیار سریع عمل می‌کند (\cref{alg:quicksort}).

\begin{latin}
\begin{algorithm}
\caption{مرتب‌سازی سریع}
\label{alg:quicksort}
\begin{algorithmic}[1]
    \Function{QuickSort}{$A$, $low$, $high$}
        \If{$low < high$}
            \State $p \gets \Call{Partition}{A, low, high}$
            \State $\Call{QuickSort}{A, low, p-1}$
            \State $\Call{QuickSort}{A, p+1, high}$
        \EndIf
    \EndFunction
\end{algorithmic}
\end{algorithm}
\end{latin}

\section{الگوریتم اقلیدس}
\label{sec:sorting-euclidean}

الگوریتم اقلیدس (\cref{alg:euclidean}) برای محاسبه‌ی
بزرگ‌ترین مقسوم‌علیه مشترک:

\begin{latin}
\begin{algorithm}
\caption{الگوریتم اقلیدس}
\label{alg:euclidean}
\begin{algorithmic}[1]
    \Require Non-negative integers $a, b$
    \Ensure $\gcd(a, b)$
    \While{$b \neq 0$}
        \State $r \gets a \bmod b$
        \State $a \gets b$
        \State $b \gets r$
    \EndWhile
    \State \Return $a$
\end{algorithmic}
\end{algorithm}
\end{latin}

\section{جدول مقایسه}
\label{sec:sorting-comparison}

جدول \cref{tab:sorting-comparison} پیچیدگی
الگوریتم‌های مرتب‌سازی را مقایسه می‌کند:

\begin{matintable}{مقایسه پیچیدگی الگوریتم‌های مرتب‌سازی}{tab:sorting-comparison}
\begin{tblr}{
    width = \textwidth,
    colspec = {l X[c] X[c] X[c] X[c]},
    hlines, vlines,
    colsep = 6pt,
    rowsep = 4pt,
    row{1} = {font=\bfseries},
}
الگوریتم & بهترین & متوسط & بدترین & فضا \\
حبابی & \lr{$O(n)$} & \lr{$O(n^2)$} & \lr{$O(n^2)$} & \lr{$O(1)$} \\
ادغامی & \lr{$O(n \log n)$} & \lr{$O(n \log n)$} & \lr{$O(n \log n)$} & \lr{$O(n)$} \\
سریع & \lr{$O(n \log n)$} & \lr{$O(n \log n)$} & \lr{$O(n^2)$} & \lr{$O(\log n)$} \\
هیپی & \lr{$O(n \log n)$} & \lr{$O(n \log n)$} & \lr{$O(n \log n)$} & \lr{$O(1)$} \\
\end{tblr}
\end{matintable}

%========================================================================
% PART 3 — ADVANCED TOPICS
%========================================================================
\part{مفاهیم پیشرفته}

%========================================================================
% CHAPTER 5 — OBJECT-ORIENTED PROGRAMMING
%========================================================================
\chapter{برنامه‌نویسی شیءگرا}
\label{chap:oop}

\section{کلاس و شیء}
\label{sec:oop-classes}

برنامه‌نویسی شیءگرا\index{برنامه‌نویسی شیءگرا}
(\lr{OOP}) روشی برای سازمان‌دهی کد بر اساس مفاهیم
«کلاس»\index{کلاس} و «شیء»\index{شیء} است.

\begin{definition}[کلاس]
    کلاس\index{کلاس!تعریف} قالبی برای ساخت اشیاء
    است که ویژگی‌ها (\lr{attributes}) و رفتارهای
    (\lr{methods}) مشترک را تعریف می‌کند.
    \label{def:class}
\end{definition}

نمونه‌ی کلاس در \cref{code:bank-account}:

\begin{latin}
\begin{minted}{python}
class BankAccount:
    """
    A simple bank account class.

    Attributes:
        account_number: Unique account identifier.
        owner: Name of the account owner.
        balance: Current account balance.
    """

    def __init__(
        self,
        account_number: str,
        owner: str,
        initial_balance: float = 0.0
    ):
        self.account_number = account_number
        self.owner = owner
        self.balance = initial_balance

    def deposit(self, amount: float) -> None:
        """Add money to the account."""
        if amount > 0:
            self.balance += amount

    def withdraw(self, amount: float) -> bool:
        """Withdraw money if sufficient balance."""
        if 0 < amount <= self.balance:
            self.balance -= amount
            return True
        return False


# Create and use accounts
account = BankAccount("IR-1234", "Ali", 1000)
account.deposit(500)
success = account.withdraw(200)
print(f"Balance: {account.balance}")  # 1300
\end{minted}
\end{latin}
\codecaption{کلاس حساب بانکی در پایتون}
\label{code:bank-account}

\section{مفاهیم اصلی OOP}
\label{sec:oop-concepts}

\begin{itemize}
    \item \textbf{کپسوله‌سازی}\index{کپسوله‌سازی}:
    پنهان‌سازی جزئیات پیاده‌سازی.
    \item \textbf{وراثت}\index{وراثت}: ایجاد کلاس
    جدید بر اساس کلاس موجود.
    \item \textbf{چندریختی}\index{چندریختی}: استفاده
    از یک رابط برای انواع مختلف.
\end{itemize}

%========================================================================
% CHAPTER 6 — DATA STRUCTURES
%========================================================================
\chapter{ساختمان داده‌های پیشرفته}
\label{chap:data-structures}

\section{درخت جستجوی دودویی}
\label{sec:ds-bst}

درخت جستجوی دودویی\index{درخت!جستجوی دودویی}
(\lr{BST}) ساختمان داده‌ای است که در آن هر گره
حداکثر دو فرزند دارد و مقدار فرزند چپ کوچک‌تر و
فرزند راست بزرگ‌تر از گره والد است
(\cref{fig:bst}).

\begin{figure}[h]
\centering
\begin{latin}
\begin{tikzpicture}[
  scale=0.9,
  transform shape,
  level distance=2cm,
  level 1/.style={sibling distance=4cm},
  level 2/.style={sibling distance=2cm}
]
  \node[treenode] {8}
    child {node[treenode] {3}
      child {node[treenode] {1}}
      child {node[treenode] {6}
        child {node[treenode] {4}}
        child {node[treenode] {7}}
      }
    }
    child {node[treenode] {10}
      child[missing]
      child {node[treenode] {14}
        child {node[treenode] {13}}
        child[missing]
      }
    };
\end{tikzpicture}
\end{latin}
\caption{درخت جستجوی دودویی}
\label{fig:bst}
\end{figure}

\section{صف و پشته}
\label{sec:ds-queue-stack}

جدول \cref{tab:queue-stack} صف و پشته را مقایسه می‌کند:

\begin{matintable}{مقایسه صف و پشته}{tab:queue-stack}
\begin{tblr}{
    width = \textwidth,
    colspec = {l X[c] X[c]},
    hlines, vlines,
    colsep = 8pt,
    rowsep = 4pt,
    row{1} = {font=\bfseries},
}
ویژگی & صف (\lr{Queue}) & پشته (\lr{Stack}) \\
اصل & \lr{FIFO}\index{FIFO} & \lr{LIFO}\index{LIFO} \\
درج & انتها (\lr{enqueue}) & بالا (\lr{push}) \\
حذف & ابتدا (\lr{dequeue}) & بالا (\lr{pop}) \\
کاربرد & صف چاپ، \lr{BFS} & بازگشت، \lr{DFS} \\
\end{tblr}
\end{matintable}

%========================================================================
% CHAPTER 7 — GRAPH ALGORITHMS
%========================================================================
\chapter{الگوریتم‌های گراف}
\label{chap:graphs}

\section{جستجوی اول سطح}
\label{sec:graphs-bfs}

جستجوی اول سطح (\cref{alg:bfs}):

\begin{latin}
\begin{algorithm}
\caption{جستجوی اول سطح (\lr{BFS})}
\label{alg:bfs}
\begin{algorithmic}[1]
    \Require Graph $G = (V, E)$, start vertex $s$
    \Ensure Shortest distances from $s$ to all vertices

    \ForAll{$v \in V$}
        \State $dist[v] \gets \infty$
    \EndFor
    \State $dist[s] \gets 0$
    \State Queue $Q \gets [s]$

    \While{$Q$ is not empty}
        \State $u \gets Q.\text{dequeue}()$
        \ForAll{neighbors $v$ of $u$}
            \If{$dist[v] = \infty$}
                \State $dist[v] \gets dist[u] + 1$
                \State $Q.\text{enqueue}(v)$
            \EndIf
        \EndFor
    \EndWhile
    \State \Return $dist$
\end{algorithmic}
\end{algorithm}
\end{latin}

\section{الگوریتم دایکسترا}
\label{sec:graphs-dijkstra}

الگوریتم دایکسترا (\cref{alg:dijkstra}):

\begin{latin}
\begin{algorithm}
\caption{الگوریتم دایکسترا}
\label{alg:dijkstra}
\begin{algorithmic}[1]
    \Require Weighted graph $G = (V, E)$, source $s$
    \Ensure Shortest distances from $s$

    \ForAll{$v \in V$}
        \State $dist[v] \gets \infty$
    \EndFor
    \State $dist[s] \gets 0$
    \State $PQ \gets$ priority queue of $(dist[v], v)$

    \While{$PQ$ is not empty}
        \State $(d, u) \gets PQ.\text{extract\_min}()$
        \If{$d > dist[u]$} \textbf{continue} \EndIf
        \ForAll{neighbors $v$ of $u$}
            \If{$dist[u] + w(u,v) < dist[v]$}
                \State $dist[v] \gets dist[u] + w(u,v)$
                \State $PQ.\text{insert}(dist[v], v)$
            \EndIf
        \EndFor
    \EndWhile
    \State \Return $dist$
\end{algorithmic}
\end{algorithm}
\end{latin}

%========================================================================
% CHAPTER 8 — EXERCISES
%========================================================================
\chapter{تمرین‌های پایانی}
\label{chap:exercises}

\section{تمرین‌های برنامه‌نویسی}
\label{sec:exercises-programming}

\begin{exercise}
    تابعی به نام \inlcode{is_prime} بنویسید که یک
    عدد صحیح دریافت کند و مشخص کند آیا آن عدد اول
    است یا خیر.
    \label{exc:is-prime}
\end{exercise}

\begin{solution}
    \begin{latin}
\begin{minted}{python}
def is_prime(n: int) -> bool:
    """Check if n is a prime number."""
    if n < 2:
        return False
    for i in range(2, int(n**0.5) + 1):
        if n % i == 0:
            return False
    return True

# Test
print(is_prime(17))  # True
print(is_prime(100)) # False
\end{minted}
    \end{latin}
\end{solution}

\begin{exercise}
    الگوریتم جستجوی دودویی (\cref{alg:binary}) را
    به‌صورت تابع پایتون پیاده‌سازی کنید.
    \label{exc:binary-search}
\end{exercise}

\begin{solution}
    \begin{latin}
\begin{minted}{python}
def binary_search(arr: list, target: int) -> int:
    """Binary search in sorted array."""
    left, right = 0, len(arr) - 1
    while left <= right:
        mid = left + (right - left) // 2
        if arr[mid] == target:
            return mid
        elif arr[mid] < target:
            left = mid + 1
        else:
            right = mid - 1
    return -1

# Test
arr = [1, 3, 5, 7, 9, 11, 13]
print(binary_search(arr, 7))  # 3
\end{minted}
    \end{latin}
\end{solution}

\section{تمرین‌های ریاضی}
\label{sec:exercises-math}

\begin{exercise}
    ثابت کنید $\sqrt{2}$ عددی گنگ است.
    \label{exc:sqrt2}
\end{exercise}

\begin{proof}
    از برهان خلف استفاده می‌کنیم. فرض کنید $\sqrt{2}$
    گویا باشد: $\sqrt{2} = \frac{a}{b}$ که
    $\gcd(a,b) = 1$.

    \[
    2 = \frac{a^{2}}{b^{2}} \implies a^{2} = 2b^{2}
    \]

    بنابراین $a^{2}$ زوج است، پس $a$ زوج است.
    $a = 2k \implies a^{2} = 4k^{2} = 2b^{2}
    \implies b^{2} = 2k^{2}$.

    بنابراین $b$ نیز زوج است. اما $\gcd(a,b) \geq 2$
    با فرض $\gcd(a,b) = 1$ تناقض دارد.
    پس $\sqrt{2}$ گنگ است.
\end{proof}

%========================================================================
% CHAPTER 9 — CONCLUSION
%========================================================================
\chapter{نتیجه‌گیری}
\label{chap:conclusion}

در این کتاب، سفری از مبانی برنامه‌نویسی تا
الگوریتم‌های پیشرفته داشتیم. از پایتون\index{پایتون}
شروع کردیم، با الگوریتم‌های
جستجو\index{الگوریتم!جستجو} و
مرتب‌سازی\index{الگوریتم!مرتب‌سازی} ادامه دادیم، و
در نهایت به ساختمان داده‌های پیشرفته\index{ساختمان داده}
و الگوریتم‌های گراف\index{الگوریتم!گراف} رسیدیم.

برای مطالعه‌ی بیشتر، منابع معرفی‌شده در بخش
کتاب‌نامه (\cite{knuth1984,cormen2009}) توصیه
می‌شوند.

\begin{flushleft}
    \textbf{پایان} \\
    مهر ۱۴۰۵
\end{flushleft}

%------------------------------------------------------------------------
% BACK MATTER (symmetric margins)
%------------------------------------------------------------------------
\backmatter
\frontmattergeometry

% Bibliography — NOT wrapped in \begin{latin}
\printbibliography[title={منابع و مراجع}]

% Index
\printindex

%------------------------------------------------------------------------
% BACK COVER (must be at the very end)
%------------------------------------------------------------------------
\authorbio{%
    تیم توسعه‌ی \lr{MatinBook}، گروهی از پژوهشگران و
    برنامه‌نویسان فارسی‌زبان است که با هدف ارائه‌ی یک
    راه‌حل حرفه‌ای برای حروف‌چینی کتاب‌های فنی فارسی
    فعالیت می‌کند. این پروژه حاصل سال‌ها تجربه در حوزه‌ی
    نشر دیجیتال و برنامه‌نویسی است.
}

\makebackcover{%
    این کتاب راهنمای کاملی برای یادگیری برنامه‌نویسی
    و الگوریتم است که با استفاده از کلاس \lr{MatinBook}
    نسخه‌ی \lr{\matinbookversion} حروف‌چینی شده است.

    \vspace{0.5cm}

    \textbf{ویژگی‌های کتاب:}
    \begin{itemize}
        \item آموزش گام‌به‌گام برنامه‌نویسی پایتون
        \item پیاده‌سازی الگوریتم‌های کلاسیک
        \item تحلیل پیچیدگی و مقایسه الگوریتم‌ها
        \item نمودارها و شبه‌کدهای استاندارد
        \item تمرین‌های متنوع با راه‌حل تشریحی
        \item منابع معتبر و نمایه‌ی جامع
    \end{itemize}

    مناسب برای دانشجویان، مدرسان و علاقه‌مندان
    به علوم کامپیوتر.
}

\end{document}